\documentclass[11pt]{article}
\usepackage{amsmath,amssymb,amsthm}
\newtheorem{lemma}{Lemma}
\title{Filtered reciprocal kernels and a failed Chebyshev shortcut}
\date{24 November 2026}
\author{Label audit research checkpoint}
\begin{document}
\maketitle

\paragraph{Status.}
The Laplace identities and reductions below are written proofs. The
counterexample is an exact finite rational interval computation with a
short continuity consequence. No statement in this note is Lean checked.
No polynomial or linear general zero bound for the product ratios is
claimed, and the requested asymptotic upper bound remains open.

\section{The exact one-variable question}

For positive weights $w_1,\ldots,w_c$ and $d>\sum_iw_i$, the full Boolean
reciprocal filter satisfies
\[
 K_w(d)=\sum_{S\subseteq[c]}\frac{(-1)^{c-|S|}}{d-\sum_{i\in S}w_i}
       =\int_0^\infty e^{-dt}\prod_i(e^{w_it}-1)\,dt.
\]
This follows by expanding the finite product and integrating each
exponential. Thus a difference of two shifted filters is the Laplace
transform of
\[
 G(t)=A e^{ut}\prod_i(e^{w_it}-1)
      -B e^{vt}\prod_i(e^{z_it}-1),\qquad A,B>0.
\]
On $t>0$ the sign changes of $G$ are the crossings of one positive
product ratio through a constant. Centering each exponential factor
reduces its logarithm to
\[
 R(t)=C+\gamma t+\sum_i\log\sinh(a_it)
                         -\sum_i\log\sinh(b_it),
 \qquad a_i,b_i>0.
\]
Its derivative is
\[
 R'(t)=\gamma+\sum_i a_i\coth(a_it)-\sum_i b_i\coth(b_it).
\]

If $G$ has at most $k$ sign changes, its Laplace transform has at most
$k$ distinct zeros on any common interval of convergence. One proof
assumes $k+1$ distinct zeros and takes a linear combination of the
corresponding $k+1$ exponential functions with zeros at the $k$ sign
changes of $G$. The exponential Chebyshev property, proved by induction
using Rolle's theorem, makes those zeros the only zeros and makes them
simple. Multiplying the combination by $-1$ if necessary, its product
with $G$ has a constant nonnegative sign and is positive on an interval.
Its integral is positive, whereas each of the assumed Laplace zeros
makes the same integral zero. This is a contradiction. The zero function
is excluded from the zero-count statement.

Consequently an $O(c)$ or $O(c^C)$ bound for the sign changes of $R$
would give a polynomial zero bound for a difference of two filtered
kernels. Such a bound would obstruct the proposed constant-family
scalar decoder with superlogarithmic word length. The reduction does
not supply the bound.

\section{The scaling family is not Chebyshev in all orders}

A tempting first lemma is that the functions
\[
 t\longmapsto a\coth(at)-\frac1t
\]
form a Chebyshev system for arbitrary distinct positive scales $a$.
That lemma is false. Put $g(x)=x\coth x-1$ and fix the ten rational
scales
\[
 a_j=1+\frac{j}{1000},\qquad 0\leq j<10.
\]
For $T>0$, form the ten by ten matrix
\[
 M(T)_{ij}=g\left(T\left(1+\frac{i}{1000}\right)
                         \left(1+\frac{j}{1000}\right)\right).
\]
The companion verifier
\texttt{JointCellCothChebyshevCertificate.py} proves by rational interval
arithmetic that
\[
 \det M(1)<0,\qquad \det M(4)>0.
\]
The approximate magnitudes, not used in verification, are respectively
$-1.56178\cdot10^{-256}$ and $1.31084\cdot10^{-228}$.

Every exponential in the verifier has a positive rational argument below
$9$. It is enclosed by its Taylor sum through degree $600$ and the
geometric upper bound on the remaining positive terms. All subsequent
arithmetic uses outward rounded rational intervals with denominator
$10^{700}$. Interval Gaussian elimination gives determinants separated
strictly from zero. The assertions use no floating point arithmetic.

By continuity there is $T_0\in(1,4)$ for which $M(T_0)$ is singular.
The ten sample points $T_0(1+i/1000)$ remain strictly increasing. A
nonzero coefficient vector in the nullspace gives a nonzero linear
combination of the ten fixed functions $g(a_jt)$ with ten distinct
positive zeros. It is not the zero function, because the nonsingular
matrix $M(1)$ already proves linear independence. Dividing by the
positive factor $t$ gives the same counterexample for
$a_j\coth(a_jt)-1/t$.

\paragraph{Scope of the failure.}
This refutes the dimension-minus-one zero bound and the direct
all-orders Chebyshev shortcut. In a product ratio the logarithmic
derivative has unit positive and negative multiplicities, whereas a
general Chebyshev linear combination has arbitrary real coefficients.
The certificate therefore does not refute an $O(c)$ bound with a larger
constant, an $O(c^C)$ bound, or a bound specialized to the product ratios
above. None of those bounds has been proved here.

\section{Distinct weight types and the real tau connection}

There is a useful elementary bound when the number of distinct positive
weights is small, even if individual weights occur many times. Suppose
the union of the two weight multisets is contained in
$\{w_1,\ldots,w_m\}$, with the $w_j$ distinct. Write their multiplicity
differences as $\kappa_j\in\mathbb Z$ and put
\[
 U(t)=C e^{\gamma t}\prod_{j=1}^m(e^{w_jt}-1)^{\kappa_j},
 \qquad C>0.
\]
The sign changes of the density difference are the crossings of
$U(t)=1$. Since every factor is positive on $t>0$, logarithmic
differentiation is legitimate there. Multiplication by the positive
product $\prod_j(e^{w_jt}-1)$ gives
\[
 \left(\frac{U'}U\right)\prod_j(e^{w_jt}-1)
 =\sum_{S\subseteq[m]}(-1)^{m-|S|}
       \left(\gamma+\sum_{j\in S}\kappa_jw_j\right)
       e^{(\sum_{j\in S}w_j)t}.
 \tag{1}
\]
After collecting coincident exponents, the right side has at most
$2^m$ exponential terms. A nonzero combination of $T$ exponentials
with distinct real exponents has at most $T-1$ distinct real zeros:
divide by the exponential of smallest exponent, differentiate, and
induct using Rolle's theorem. Thus, unless $U$ is constant, $U'/U$
has at most $2^m-1$ zeros and $\log U$ has at most $2^m$ zeros.
If $U$ is constant, the density difference is either identically zero
or has a fixed sign. The preceding Laplace argument consequently gives
at most $2^m$ distinct zeros for every nonzero difference of the two
filtered kernels. This is a written proof, not a new Lean declaration.

The bound depends on distinct weight types, not filter order. In
particular, if each of $a$ scalar bit readers is a nonzero difference
of two such filters, using at most $m$ weight types, an ordered token
alphabet with all $2^a$ distinct strict words must satisfy
\[
 2^a\leq1+a2^m.
\]
Each transition between consecutive distinct words crosses a zero of
at least one reader, and the total number of available zeros is at most
$a2^m$. This argument does not apply to a general sum of more than two
filters without a further zero bound.

If the weights are commensurable, $w_j=\lambda k_j$ with positive
integers $k_j\leq d$, the right side of (1), as a polynomial in
$x=e^{\lambda t}$, has degree at most $\sum_j k_j\leq md$.
The improved kernel zero bound is then $md+1$, regardless of the real
shift difference $\gamma$. Merely saying that the weights are rational
does not make this a polynomial bound in $c$: clearing denominators
can make $d$ arbitrarily large.

For integer or rational weights, the density itself becomes, after
$x=e^{-\lambda t}$ and removal of a nonzero monomial factor, a sum of
two products of binomials. This lies in the class studied by
Grenet, Koiran, Portier and Strozecki, \emph{The Limited Power of
Powering}, Theorem 1, and Koiran, Portier and Tavenas,
\emph{A Wronskian Approach to the Real $\tau$-Conjecture},
Theorems 12 and 13. In their notation the number of summands is $k=2$,
the sparsity is $t=2$, and the number of distinct sparse base
polynomials is as large as $m=2c+1$. Theorem 12 gives
$t^{O(mk^2)}$, which is exponential in $c$ here, despite fixed $k$.
The bounded-base-degree result in Theorem 13 is polynomial only when
those degrees are controlled. These results do not supply the desired
polynomial bound for arbitrarily many independently chosen weights.

\paragraph{Primary sources.}
The checked source statements are at
\texttt{https://arxiv.org/pdf/1205.1015}, Theorems 12 and 13,
and \texttt{https://www-verimag.imag.fr/\string~grenetb/publis/GKPS11.pdf},
Theorem 1. The broad polynomial root-count assertion for sums of
products of sparse polynomials is the real tau-conjecture. We have not
proved that our restricted two-product family is equivalent to that
conjecture, nor that its special arbitrary-weight case is known to be
open. The point of the comparison is the exact limitation of the cited
theorems, not an unsupported impossibility claim.

\end{document}
